\section{Interim Report}

The purpose of the interim report is to demonstrate the actual progress of the project, verify that the chosen direction is appropriate, and allow adjustments before the final submission.

By the indicated date, each group must submit a brief report of 3--5 pages via Aula Virtual. The report must include at least the following elements:

\begin{itemize}
    \item \textbf{Updated problem statement and objectives.} Revise, if necessary, the initial formulation of the problem and objectives based on what has been learned so far.

    \item \textbf{Description of the proposed solution.} Explain in greater detail the chosen approach, including models, algorithms, architecture, procedures, and other relevant aspects. Diagrams or pseudocode should be included where useful.

    \item \textbf{Implementation plan.} Summarize what has been implemented so far and what remains to be done, including modules, components, and main tasks.

    \item \textbf{Dataset and preprocessing.} Describe the data effectively used so far, how it was obtained, and which preprocessing steps have been applied or are planned.

    \item \textbf{Preliminary experiments, if applicable.} Report preliminary experiments, even if they are simple or incomplete, together with the metrics used.

    \item \textbf{Risks and adjustments to the plan.} Indicate difficulties encountered, identified risks, such as data quality, time constraints, or technical complexity, and explain how the work plan will be adapted to reach the final submission.

    \item \textbf{Updated bibliography.} Provide a list of relevant papers and resources guiding the project up to this point. This list may be expanded later in the final submission.
\end{itemize}